\documentclass[11pt,a4paper]{article}

%================================================
% ENCODAGE ET LANGUE
%================================================
\usepackage[utf8]{inputenc}
\usepackage[T1]{fontenc}
\usepackage[french]{babel}

%================================================
% MISE EN PAGE
%================================================
\usepackage{geometry}
\geometry{margin=1.55cm}

%================================================
% PACKAGE GTOLI
%================================================
\usepackage{../styles/gtoli}

%================================================
% INFORMATIONS DU DOCUMENT
%================================================
\title{\textbf{Économie}\\[0.2cm]
\large Titre du chapitre}

\author{GToli}
\date{\today}

%================================================
% DOCUMENT
%================================================
\begin{document}

\maketitle

%================================================
\section{Notion économique concernée}
%================================================

\begin{cadreobjectif}
Dans ce chapitre, l’élève doit être capable de :
\begin{itemize}
    \item identifier la notion économique étudiée ;
    \item utiliser correctement le vocabulaire économique ;
    \item lire et interpréter un graphique ou un tableau ;
    \item distinguer causes, mécanismes et conséquences ;
    \item mobiliser un raisonnement économique structuré ;
    \item rédiger une réponse argumentée.
\end{itemize}
\end{cadreobjectif}

%================================================
\section{Situation de départ}
%================================================

\begin{cadrebleu}{Problématique économique}
Cette section permet d’introduire une situation économique concrète.

\medskip

On précisera :
\begin{itemize}
    \item le phénomène observé ;
    \item les acteurs concernés ;
    \item le contexte économique ;
    \item la question principale du chapitre.
\end{itemize}
\end{cadrebleu}

%================================================
\section{Vocabulaire économique}
%================================================

\begin{cadretheorie}
\begin{itemize}
    \item \textbf{Concept 1 :} définition.
    \item \textbf{Concept 2 :} définition.
    \item \textbf{Concept 3 :} définition.
    \item \textbf{Concept 4 :} définition.
\end{itemize}
\end{cadretheorie}

%================================================
\section{Rappel théorique}
%================================================

\begin{cadretheorie}
Cette section contient :
\begin{itemize}
    \item les définitions essentielles ;
    \item les mécanismes économiques ;
    \item les relations entre variables ;
    \item les indicateurs ;
    \item les limites d’interprétation ;
    \item les exemples contextualisés.
\end{itemize}

\medskip

\textbf{Point essentiel :} en économie, il faut toujours distinguer le constat, le mécanisme explicatif et la conséquence.
\end{cadretheorie}

%================================================
\section{Acteurs, variables et mécanismes}
%================================================

\begin{cadregris}{Tableau d’analyse}
\renewcommand{\arraystretch}{1.35}
\begin{tabularx}{\textwidth}{>{\bfseries}p{3.2cm}X X}
\toprule
Élément & Rôle économique & Effet possible \\
\midrule
Ménages & Consommation, épargne, travail & Effet sur la demande \\
Entreprises & Production, investissement, emploi & Effet sur l’offre \\
État & Régulation, redistribution, fiscalité & Effet sur l’activité \\
Banques & Crédit, financement & Effet sur l’investissement \\
\bottomrule
\end{tabularx}
\end{cadregris}

%================================================
\section{Graphique économique}
%================================================

\begin{cadregris}{Lecture graphique}
Insérer ici un graphique économique : offre-demande, coût, productivité, inflation, chômage, PIB, courbe de Lorenz, etc.

\medskip

\begin{center}
\begin{tikzpicture}[scale=0.9]
\draw[->,thick] (0,0) -- (6,0) node[right] {Quantité};
\draw[->,thick] (0,0) -- (0,4) node[above] {Prix};
\draw[thick] (0.8,3.5) -- (5.3,0.6) node[right] {Demande};
\draw[thick] (0.8,0.6) -- (5.3,3.5) node[right] {Offre};
\fill (3,2.05) circle (2pt);
\node[right] at (3,2.05) {Équilibre};
\end{tikzpicture}
\end{center}
\end{cadregris}

%================================================
\section{Méthode d’analyse économique}
%================================================

\begin{cadremethode}
Pour analyser une situation économique :

\begin{enumerate}
    \item identifier le phénomène économique ;
    \item repérer les acteurs concernés ;
    \item relever les variables importantes ;
    \item distinguer les causes et les conséquences ;
    \item expliquer le mécanisme ;
    \item utiliser les données ou graphiques disponibles ;
    \item nuancer l’interprétation ;
    \item rédiger une conclusion argumentée.
\end{enumerate}
\end{cadremethode}

%================================================
\section{Lecture de document}
%================================================

\begin{cadreenonce}
\textbf{Document :}

Insérer ici un article, un tableau statistique, un graphique, un extrait institutionnel ou une situation économique.

\medskip

\textbf{Question :} analyser le document en mobilisant les notions du chapitre.
\end{cadreenonce}

\begin{cadreaide}
Pour analyser un document économique :
\begin{itemize}
    \item identifier la source ;
    \item repérer les données importantes ;
    \item définir les concepts utilisés ;
    \item expliquer les relations entre variables ;
    \item formuler une conclusion nuancée.
\end{itemize}
\end{cadreaide}

%================================================
\section{Exemple modèle}
%================================================

\begin{cadreexemple}
\textbf{Situation :}

Insérer ici une situation économique représentative du chapitre.

\medskip

\textbf{Analyse :}
\begin{itemize}
    \item Quel est le phénomène étudié ?
    \item Quels sont les acteurs concernés ?
    \item Quelles sont les variables en jeu ?
    \item Quel mécanisme permet d’expliquer la situation ?
\end{itemize}

\medskip

\textbf{Réponse structurée :}

Présenter ici une réponse économique complète.
\end{cadreexemple}

%================================================
\section{Exercice d’application}
%================================================

\begin{cadreenonce}
\textbf{Exercice 1 :}

Insérer ici un exercice d’application directe.
\end{cadreenonce}

%================================================
\section{Solution détaillée}
%================================================

\begin{cadresolution}
La solution doit comprendre :
\begin{enumerate}
    \item l’identification du phénomène ;
    \item la définition des concepts ;
    \item l’analyse des données ;
    \item l’explication du mécanisme ;
    \item la conclusion argumentée.
\end{enumerate}
\end{cadresolution}

%================================================
\section{Erreurs fréquentes}
%================================================

\begin{cadreattention}
Les erreurs fréquentes en économie sont :
\begin{itemize}
    \item confondre corrélation et causalité ;
    \item décrire un graphique sans l’interpréter ;
    \item oublier les acteurs économiques ;
    \item utiliser un concept sans le définir ;
    \item tirer une conclusion trop générale ;
    \item négliger le contexte.
\end{itemize}
\end{cadreattention}

%================================================
\section{Exercices progressifs}
%================================================

\begin{cadreactivite}
\begin{enumerate}
    \item définir les concepts ;
    \item compléter un schéma économique ;
    \item lire un graphique ;
    \item interpréter un tableau de données ;
    \item analyser un document ;
    \item rédiger une réponse argumentée.
\end{enumerate}
\end{cadreactivite}

%================================================
\section{Synthèse}
%================================================

\begin{cadresynthese}
À retenir :

\begin{itemize}
    \item notion économique principale : \dotfill
    \item acteurs concernés : \dotfill
    \item variables importantes : \dotfill
    \item mécanisme central : \dotfill
    \item erreur à éviter : \dotfill
    \item conclusion économique : \dotfill
\end{itemize}
\end{cadresynthese}

%================================================
\section{Auto-évaluation}
%================================================

\begin{cadregris}{Je vérifie ma maîtrise}
\begin{itemize}
    \item[$\square$] Je connais le vocabulaire économique.
    \item[$\square$] Je sais identifier les acteurs.
    \item[$\square$] Je sais lire un graphique.
    \item[$\square$] Je sais expliquer un mécanisme.
    \item[$\square$] Je sais analyser un document.
    \item[$\square$] Je sais rédiger une réponse argumentée.
\end{itemize}
\end{cadregris}

\end{document}